\chapter{Compressed Sensing}
\label{ch:compressed-sensing}
\fancyhead[R]{\small\textit{Chapter 12: Compressed Sensing}}
\section{Overview: From Redundancy to Sparsity}
\label{sec:ch12-overview}

Lab Exercise~5.5 of Chapter~5 ended with the following observation: given
the frame $\{h_1,h_2,h_3\}$ in $\R^2$, there is not a unique coefficient
vector $c\in\R^3$ reconstructing a given signal $x$. Rather, there is an
entire affine family $c' = c + t\,k$ (where $k$ spans $\ker H$), and
different choices of $t$ produce different coefficient vectors, all
correctly representing $x$. The exercise asked you to minimize $\norm{c'}_1$
over this family and observe that the minimizer is \emph{sparse} --- it has
an exact zero entry. The comment at the end read: ``this phenomenon ---
$\ell^1$ minimization promoting exact sparsity --- is the starting point
of compressed sensing, which we will touch on again in Chapter~12.''

This is that chapter. Compressed sensing is a framework for acquiring
signals using \emph{far fewer measurements than the ambient dimension}
by exploiting the prior knowledge that the signal is sparse.
Frames are central to the theory in two ways: as the \emph{sensing
matrices} (the $n\times m$ matrices $A$ that measure a sparse signal
$x\in\R^m$ via $y = Ax\in\R^n$, with $n\ll m$), and as the
\emph{dictionaries} in which sparse representations are sought. The
frame-theoretic tools we have developed --- coherence (Chapter~9),
frame bounds (Chapter~3), and the $\ker F$ parametrization of redundant
representations (Chapter~5) --- appear throughout.

This chapter is an introduction rather than a complete treatment. Its
goal is to show precisely how the frame-theoretic ideas of this course
are the right language for stating and understanding compressed sensing,
and to give you enough machinery to implement and test sparse recovery
algorithms numerically.

\section{Sparsity and Redundant Representations}
\label{sec:sparsity}

\begin{definition}[$s$-sparse vector]
\label{def:sparse}
A vector $c\in\F^m$ is \emph{$s$-sparse} if at most $s$ of its entries
are nonzero: $|\operatorname{supp}(c)| \le s$, where $\operatorname{supp}(c)
:= \setdef{i}{c_i \ne 0}$.
\end{definition}

\begin{definition}[$\ell^1$ norm and $\ell^0$ ``norm'']
\label{def:l1-norm}
The \emph{$\ell^1$ norm} of $c = (c_1,\ldots,c_m)\in\F^m$ is
$\norm{c}_1 := \sum_{i=1}^m |c_i|$. The \emph{$\ell^0$ count}
$\norm{c}_0 := |\operatorname{supp}(c)|$ counts the nonzero entries
(it is not a norm, but the notation is standard).
\end{definition}

\subsection{Sparsity in the Frame Setting}
\label{subsec:sparsity-in-frames}

Recall from Chapter~5 (Theorem~5.4.1) that when a frame
$\{f_i\}_{i=1}^m$ has more vectors than the dimension ($m>n$),
every signal $x\in\F^n$ has infinitely many coefficient representations:
all vectors in the affine family $c + \ker F$, where $c = F^*x$ is
the analysis operator's output. Among all these representations, we
can ask: which one is \emph{sparsest}?

\begin{definition}[Sparse representation problem]
\label{def:sparse-rep}
Given a frame $\{f_i\}_{i=1}^m$ (synthesis matrix $F$) and a signal
$x\in\F^n$, the \emph{sparse representation problem} is
\begin{equation}
\label{eq:l0-problem}
  \minimize_{c\in\F^m} \norm{c}_0 \quad\text{subject to}\quad Fc = x.
\end{equation}
\end{definition}

Problem~\eqref{eq:l0-problem} is combinatorially hard in general:
finding the sparsest solution requires, in the worst case, checking all
$\binom{m}{s}$ subsets of size $s$ for $s = 1, 2, \ldots$ until one
spanning $x$ is found. The key insight of compressed sensing is that
under mild conditions, replacing the $\ell^0$ count with the $\ell^1$
norm gives the same answer, and $\ell^1$ minimization is a
\emph{convex optimization problem} solvable efficiently.

\begin{definition}[Basis pursuit]
\label{def:basis-pursuit}
The \emph{basis pursuit} problem is
\begin{equation}
\label{eq:basis-pursuit}
  \minimize_{c\in\F^m} \norm{c}_1 \quad\text{subject to}\quad Fc = x.
\end{equation}
\end{definition}

\begin{remark*}
Basis pursuit is a convex program: $\norm{c}_1$ is a convex function
and $\{c : Fc = x\}$ is an affine (hence convex) constraint set.
Convex programs are solvable in polynomial time (e.g.\ via interior-point
methods), and basis pursuit can be recast as a linear program that any
standard LP solver handles efficiently. In NumPy/SciPy, it can be
solved via \texttt{scipy.optimize.linprog} or the more convenient
\texttt{cvxpy} library. We show both approaches in
Section~\ref{sec:numpy-ch12}.
\end{remark*}

\begin{examplebox}[title={Example 12.1: The triangular frame --- $\ell^1$ finds the sparsest expansion}]
\label{ex:triangular-sparsity}
Consider the triangular frame $f_1=(1,0)^T$, $f_2=(-\tfrac12,\tfrac{\sqrt3}{2})^T$,
$f_3=(-\tfrac12,-\tfrac{\sqrt3}{2})^T$ and $x = f_1 = (1,0)^T$. The
canonical coefficient vector is $c = \tfrac{2}{3}(1,-\tfrac12,-\tfrac12)^T
= (\tfrac23,-\tfrac13,-\tfrac13)^T$ (with no zeros). The kernel of $F$
is spanned by $k = (1,1,1)^T/\sqrt3$, so all representations are
\[
  c(t) = c + tk =
  \Bigl(\tfrac23 + \tfrac{t}{\sqrt3},\;
        -\tfrac13 + \tfrac{t}{\sqrt3},\;
        -\tfrac13 + \tfrac{t}{\sqrt3}\Bigr)^T.
\]
The function $\norm{c(t)}_1$ is piecewise linear and is minimized at
$t^* = 1/\sqrt3$, giving
\[
  c^* = c(t^*) = (1,\,0,\,0)^T.
\]
This is the unique 1-sparse solution (the sparsest possible), correctly
identifying $x = f_1$ as a scalar multiple of the first frame vector.
Its $\ell^1$ norm is $1$, compared to the canonical vector's
$\norm{c}_1 = \tfrac23+\tfrac13+\tfrac13=\tfrac43$.

The proof is a direct calculation: with $t^*=1/\sqrt3$, the second and
third entries both equal $-\tfrac13+\tfrac{1}{\sqrt3}\cdot\tfrac{1}{\sqrt3}
= -\tfrac13+\tfrac13 = 0$, so $c^* = (1,0,0)^T$ exactly.
\end{examplebox}

\begin{keyidea}
\textbf{$\ell^1$ minimization over $\ker F$ finds the sparsest
representation of $x$ in the frame.} The canonical coefficients
$c=F^*\tilde x$ (from the analysis operator) minimize
$\norm{c}_2$ (Chapter~5, Theorem~5.5.1), but not $\norm{c}_1$.
$\ell^1$ minimization --- basis pursuit --- solves a different problem
and generally finds a sparser answer. These are two orthogonal
objectives, and which one you want depends on your application.
\end{keyidea}

\section{Compressed Sensing: The Measurement Model}
\label{sec:cs-model}

Basis pursuit finds sparse expansions in a given frame. Compressed
sensing goes further: it acquires a signal $x\in\R^m$ using far fewer
than $m$ measurements, and recovers it exactly by exploiting sparsity.

\begin{definition}[Compressed sensing model]
\label{def:cs-model}
Let $A\in\R^{n\times m}$ with $n\ll m$ be the \emph{sensing matrix}
(or \emph{measurement matrix}). The receiver observes $y = Ax\in\R^n$
(noiseless) or $y = Ax + e\in\R^n$ (noisy, $\norm{e}_2\le\eta$).
Given $y$, $A$, and the prior that $x$ is $s$-sparse ($\norm{x}_0\le s$),
recover $x$.
\end{definition}

Without the sparsity prior, $y=Ax$ with $n<m$ is underdetermined and
has infinitely many solutions. With the sparsity prior, the system
can be uniquely solvable --- but finding the sparsest solution is
$\ell^0$-hard. Basis pursuit gives the tractable surrogate:
\begin{equation}
\label{eq:cs-bp}
  \hat x = \argmin_{z\in\R^m}\norm{z}_1 \quad
  \text{subject to}\quad Az = y.
\end{equation}

\begin{remark*}
The columns of $A$ form a frame for $\R^n$ (assuming $A$ has rank $n$,
which is necessary for any reconstruction to be possible). Here the
frame is $\{a_i\}_{i=1}^m$ where $a_i$ is the $i$-th column of $A$,
and the measurement $y_j = \langle x, a_j^{\text{row}}\rangle$ is the
inner product of the signal $x$ with the $j$-th \emph{row} of $A$.
Both perspectives --- columns as frame vectors for $\R^n$, rows as
measurement functionals on $\R^m$ --- appear in the literature; we use
the row perspective throughout.
\end{remark*}

\section{When Does Basis Pursuit Recover Sparse Signals?}
\label{sec:recovery-conditions}

The central question is: for which matrices $A$ does basis pursuit
\eqref{eq:cs-bp} recover every $s$-sparse signal $x$ exactly?
Two complementary answers are the \emph{coherence condition} and the
\emph{restricted isometry property}.

\subsection{The Coherence Condition}
\label{subsec:coherence-condition}

Recall coherence from Chapter~9 (Definition~9.2.1): the coherence of
a unit-column-norm matrix $A$ is $\mu(A) := \max_{i\ne j}
|\inner{a_i}{a_j}|$. Small coherence means the columns of $A$ are
``maximally incoherent'' --- no two look alike.

\begin{theorem}[Coherence condition for exact recovery]
\label{thm:coherence-recovery}
Let $A\in\R^{n\times m}$ have unit-norm columns and coherence $\mu$.
If $x\in\R^m$ is $s$-sparse with
\begin{equation}
\label{eq:coherence-condition}
  s < \frac{1}{2}\!\left(1 + \frac{1}{\mu}\right),
\end{equation}
then basis pursuit \eqref{eq:cs-bp} with $y=Ax$ recovers $x$ exactly.
\end{theorem}

\begin{proof}[Proof sketch]
We show that if $s < (1+1/\mu)/2$ then $x$ is the \emph{unique}
$s$-sparse solution to $Az=y$, hence the $\ell^0$ problem has a unique
solution; a separate argument (which we omit) shows that in this case
basis pursuit gives the same answer as the $\ell^0$ problem.

Suppose $Az=y=Ax$ with $z\ne x$ and $\norm{z}_0\le s$. Then $A(z-x)=0$
with $h := z-x \ne 0$ and $\norm{h}_0 \le 2s$. Write
$h=\sum_{i\in T}h_i e_i$ where $|T|\le 2s$. Then $0=Ah=\sum_{i\in T}
h_i a_i$. Taking the inner product with $a_j$ for $j\in T$:
\[
  \norm{a_j}^2 |h_j|
  = \left|\sum_{i\in T,i\ne j} h_i\inner{a_i}{a_j}\right|
  \le \mu \sum_{i\in T,i\ne j}|h_i|
  \le \mu(2s-1)\norm{h}_\infty,
\]
since $|\inner{a_i}{a_j}|\le\mu$ and there are at most $2s-1$ terms
in the sum. Since $\norm{a_j}^2=1$ and this holds for all $j\in T$,
we get $\norm{h}_\infty \le \mu(2s-1)\norm{h}_\infty$, which requires
$\mu(2s-1)\ge1$, i.e.\ $s\ge(1+1/\mu)/2$. This contradicts our
assumption $s<(1+1/\mu)/2$. Hence no such $z\ne x$ exists, and the
$s$-sparse solution is unique.
\end{proof}

\begin{remark*}
Condition~\eqref{eq:coherence-condition} is sufficient but not
necessary: basis pursuit can and does recover signals even when
$s\ge(1+1/\mu)/2$ in practice. The bound is pessimistic because
it assumes the worst-case alignment of $s$ columns. Empirically,
random matrices with $n = O(s\log m)$ measurements recover $s$-sparse
signals with high probability even when $s$ is much larger than the
coherence bound allows.
\end{remark*}

\begin{examplebox}[title={Example 12.2: Coherence bound for the triangular frame}]
\label{ex:coherence-bound-triangle}
The triangular frame has coherence $\mu = \tfrac12$ (all off-diagonal
inner products equal $-\tfrac12$, so $|\inner{f_i}{f_j}|=\tfrac12$
for $i\ne j$; this is the Welch bound for $(n,m)=(2,3)$, Chapter~9).
The coherence condition guarantees recovery of $s$-sparse signals for
\[
  s < \tfrac12\bigl(1+\tfrac1{1/2}\bigr)
  = \tfrac12(1+2) = \tfrac32,
\]
i.e.\ $s\le 1$. Since the ambient dimension is $n=2$, this means basis
pursuit recovers any $1$-sparse signal --- any multiple of a single
frame vector --- exactly. Example~\ref{ex:triangular-sparsity}
confirmed this directly for $x=f_1$.
\end{examplebox}

\subsection{The Restricted Isometry Property}
\label{subsec:rip}

The coherence condition is simple to state and check, but it gives
pessimistic bounds that preclude recovery for larger sparsity levels.
The \emph{restricted isometry property} (RIP) captures a more refined
condition that allows recovery of much sparser signals.

\begin{definition}[Restricted isometry property, RIP]
\label{def:rip}
A matrix $A\in\R^{n\times m}$ satisfies the \emph{RIP of order $s$}
with constant $\delta_s\in(0,1)$ if, for every $s$-sparse vector
$z\in\R^m$,
\begin{equation}
\label{eq:rip}
  (1-\delta_s)\norm{z}_2^2 \;\le\; \norm{Az}_2^2 \;\le\; (1+\delta_s)\norm{z}_2^2.
\end{equation}
The smallest valid $\delta_s$ is called the \emph{RIP constant} of
order $s$.
\end{definition}

\begin{remark*}
Compare~\eqref{eq:rip} to the frame inequality (Definition~2.1.1):
$A\norm{z}^2\le\norm{F^*z}^2\le B\norm{z}^2$ for all $z$. The RIP is
a \emph{restricted} frame inequality: it asks for the frame condition
to hold, but only for $s$-sparse vectors $z$ (not for all $z$). The
RIP constant $\delta_s$ is the worst-case deviation from an isometry
over all $s$-sparse vectors --- it measures how far $A$ is from
being a perfect isometry on sparse signals. If $\delta_s=0$, then
$A$ maps every $s$-sparse signal to a vector of the same norm: perfect
preservation. Larger $\delta_s$ means more distortion.

In frame language: $A$ satisfies RIP-$s$ with constant $\delta_s$ if
and only if \emph{every} $n\times s$ submatrix of $A$ (formed by
taking $s$ columns) has frame bounds $A_{\min}\ge 1-\delta_s$ and
$B_{\max}\le 1+\delta_s$. This is exactly the condition that every
$s$-column subframe of $A$ is a well-conditioned frame for $\R^n$.
\end{remark*}

\begin{theorem}[RIP guarantees exact recovery]
\label{thm:rip-recovery}
If $A$ satisfies the RIP of order $2s$ with constant $\delta_{2s} <
\sqrt{2}-1 \approx 0.414$, then for every $s$-sparse $x\in\R^m$,
basis pursuit with $y=Ax$ recovers $x$ exactly.
\end{theorem}

We state this theorem without proof, as the argument is substantially
more involved than the coherence condition proof and requires tools
beyond the scope of this course. The key point is that the threshold
$\delta_{2s}<\sqrt{2}-1$ is sharp: there exist matrices with
$\delta_{2s}$ slightly above this threshold for which basis pursuit
fails on some $s$-sparse signal.

\begin{proposition}[Coherence and RIP]
\label{prop:coherence-rip}
If $A$ has unit-norm columns and coherence $\mu$, then
\[
  \delta_s \;\le\; (s-1)\mu
\]
for all $s\ge 1$. In particular, if $\mu < (\sqrt{2}-1)/(2s-1)$,
then $\delta_{2s}\le(2s-1)\mu<\sqrt{2}-1$, so the RIP condition for
exact recovery is satisfied.
\end{proposition}

\begin{proof}
Let $z\in\R^m$ be $s$-sparse with support $T$, $|T|=s$. Then
\[
  \norm{Az}^2 = \sum_{i\in T}\sum_{j\in T}z_iz_j\inner{a_i}{a_j}
  = \norm{z}^2 + \sum_{i\ne j,\,i,j\in T}z_iz_j\inner{a_i}{a_j}.
\]
The cross terms satisfy $|z_iz_j\inner{a_i}{a_j}|\le\mu|z_iz_j|$, so
\[
  \bigl|\norm{Az}^2 - \norm{z}^2\bigr|
  \le \mu\sum_{i\ne j\in T}|z_iz_j|
  \le \mu(s-1)\norm{z}^2,
\]
where the last step uses $\sum_{i\ne j}|z_iz_j|\le (s-1)\sum_i z_i^2
= (s-1)\norm{z}^2$ (a consequence of the AM-GM or Cauchy-Schwarz
inequality). Hence $\delta_s\le(s-1)\mu$.
\end{proof}

\begin{remark*}
Proposition~\ref{prop:coherence-rip} shows that coherence controls
the RIP constant, but gives a weaker bound than direct RIP analysis.
Random matrices --- Gaussian, Bernoulli, or random partial Fourier ---
achieve RIP constants $\delta_s \lesssim \sqrt{s/n}$ with high
probability when $n = O(s\log(m/s))$ measurements are taken. This is
far better than coherence-based bounds, which require $n = O(s^2\mu^{-2})$.
\end{remark*}

\section{The Frame-Theoretic View of Sensing}
\label{sec:frame-view}

Compressed sensing and frame theory are two perspectives on the same
algebraic structure. This section makes the connections explicit.

\begin{proposition}[The sensing matrix is a frame]
\label{prop:sensing-frame}
A sensing matrix $A\in\R^{n\times m}$ with $\rank(A)=n$ and columns
$\{a_i\}_{i=1}^m$ forms a frame for $\R^n$. The RIP constant $\delta_s$
is the worst-case deviation of $s$-column sub-frame bounds from $(1,1)$:
\[
  \delta_s = \max_{|T|=s}\max\!\bigl(1-\lambda_{\min}(A_T^TA_T),\;
  \lambda_{\max}(A_T^TA_T)-1\bigr),
\]
where $A_T$ denotes the submatrix of $A$ with columns indexed by $T$.
\end{proposition}

\begin{proof}
$A$ has $\rank n$ by assumption, so its columns span $\R^n$ and
Theorem~2.2.1 gives a frame. The RIP formula is immediate from
Definition~\ref{def:rip}: $\norm{Az}^2 = \norm{A_Tc_T}^2 = c_T^T A_T^T
A_T c_T$ where $c_T=(z_i)_{i\in T}$, and the frame bounds of
$\{a_i\}_{i\in T}$ are $\lambda_{\min}(A_T^TA_T)$ and
$\lambda_{\max}(A_T^TA_T)$ (by Theorem~3.5.1 applied to the sub-frame).
\end{proof}

\begin{keyidea}
\textbf{The RIP is the frame condition, restricted to $s$-column
submatrices.} A sensing matrix $A$ satisfies RIP-$s$ if and only if
every $s$-column subframe $\{a_i\}_{i\in T}$ is a
$(1-\delta_s,1+\delta_s)$-frame for $\R^n$. Small $\delta_s$ means
all $s$-column subframes are nearly tight: every $s$ columns span
$\R^n$ approximately isometrically. The closer $\delta_s$ is to zero,
the better the measurement matrix preserves the geometry of sparse
signals.
\end{keyidea}

\begin{proposition}[ETFs are optimal sensing matrices]
\label{prop:etf-optimal-sensing}
Among all $n\times m$ matrices with unit-norm columns, ETFs (equiangular
tight frames, Chapter~9) have the smallest possible worst-case coherence
$\mu = \sqrt{(m-n)/(n(m-1))}$ (the Welch bound). Since
Proposition~\ref{prop:coherence-rip} bounds $\delta_s\le(s-1)\mu$, ETFs
give the best coherence-based RIP bounds among all unit-norm sensing
matrices of the same size. They are therefore optimal sensing matrices
in the coherence-minimization sense.
\end{proposition}

\section{NumPy Implementation}
\label{sec:numpy-ch12}

\begin{lstlisting}[caption={Basis pursuit and compressed sensing in Python}]
import numpy as np
from scipy.optimize import linprog, minimize_scalar
from scipy.linalg import null_space

# -------------------------------------------------------------------
# Example 12.1: L1 minimization in the triangular frame
# -------------------------------------------------------------------
f1 = np.array([1.0, 0.0])
f2 = np.array([-0.5, np.sqrt(3)/2])
f3 = np.array([-0.5, -np.sqrt(3)/2])
F = np.column_stack([f1, f2, f3])   # synthesis matrix (2 x 3)

x = f1   # signal to represent: x = (1, 0)^T

# Canonical coefficients via the frame operator
S = F @ F.T
c_can = F.T @ np.linalg.solve(S, x)
print(f"Canonical c = {c_can.round(4)}")   # (2/3, -1/3, -1/3)
print(f"  ||c||_1 = {np.sum(np.abs(c_can)):.4f}")  # 4/3

# L1 minimization over the affine family c + t*k
k = null_space(F)[:, 0]   # unit kernel vector (1,1,1)/sqrt(3)

def l1_objective(t):
    return np.sum(np.abs(c_can + t * k))

result = minimize_scalar(l1_objective, bounds=(-3, 3), method='bounded')
t_opt = result.x
c_l1 = c_can + t_opt * k
print(f"L1-optimal c = {c_l1.round(6)}")   # (1, 0, 0) to machine precision
print(f"  ||c||_1 = {np.sum(np.abs(c_l1)):.4f}")   # 1.0
print(f"  Reconstruction ok: {np.allclose(F @ c_l1, x)}")

# -------------------------------------------------------------------
# Compressed sensing: basis pursuit via linear programming
# -------------------------------------------------------------------
def basis_pursuit(A, y):
    """
    Solve min ||c||_1 s.t. Ac = y via linear programming.

    Reformulation: min 1^T t
                   s.t. Ac = y,  -t <= c <= t,  t >= 0
    Variables: [c (m,), t (m,)]
    """
    n, m = A.shape
    # Objective: minimize sum of t_i
    c_obj = np.zeros(2 * m); c_obj[m:] = 1.0
    # Equality: A c = y
    A_eq = np.hstack([A, np.zeros((n, m))])
    b_eq = y
    # Inequality: c_i - t_i <= 0  and  -c_i - t_i <= 0
    A_ub = np.block([[ np.eye(m), -np.eye(m)],
                     [-np.eye(m), -np.eye(m)]])
    b_ub = np.zeros(2 * m)
    # Bounds: c_i free, t_i >= 0
    bounds = [(None, None)] * m + [(0, None)] * m
    res = linprog(c_obj, A_ub=A_ub, b_ub=b_ub,
                  A_eq=A_eq, b_eq=b_eq, bounds=bounds, method='highs')
    return res.x[:m]

def coherence(A):
    """Coherence of matrix A with unit-norm columns."""
    G = A.T @ A
    off = np.abs(G - np.diag(np.diag(G)))
    return off.max()

def rip_constant(A, s):
    """
    Estimate RIP-s constant: worst-case submatrix frame-bound deviation.
    Samples up to 2000 random s-subsets for efficiency.
    """
    from itertools import combinations
    import random
    n, m = A.shape
    all_subsets = list(combinations(range(m), s))
    if len(all_subsets) > 2000:
        all_subsets = random.sample(all_subsets, 2000)
    delta = 0.0
    for T in all_subsets:
        A_T = A[:, list(T)]
        eigs = np.linalg.eigvalsh(A_T.T @ A_T)
        delta = max(delta, abs(eigs[-1] - 1), abs(1 - eigs[0]))
    return delta

# -------------------------------------------------------------------
# Demo: recover a 2-sparse signal in R^40 from 15 measurements
# -------------------------------------------------------------------
rng = np.random.default_rng(42)
n_meas, m_sig, s_spar = 15, 40, 2

# Sensing matrix: random Gaussian with normalized columns
A = rng.standard_normal((n_meas, m_sig))
A = A / np.linalg.norm(A, axis=0)

mu_A = coherence(A)
delta_2s = rip_constant(A, 2 * s_spar)
print(f"\nSensing matrix ({n_meas}x{m_sig}): mu={mu_A:.4f}, "
      f"estimated delta_{{2s}}={delta_2s:.4f}")
print(f"Coherence bound: s < {0.5*(1 + 1/mu_A):.2f}")
print(f"RIP recovery condition: delta_{{2s}} < {np.sqrt(2)-1:.4f}, "
      f"holds: {delta_2s < np.sqrt(2)-1}")
# Note: the RIP condition is sufficient, not necessary.
# Basis pursuit can succeed even when delta_{2s} >= sqrt(2)-1;
# the condition is a worst-case guarantee over ALL 2s-sparse signals.

# True sparse signal
x_true = np.zeros(m_sig)
x_true[5] = 3.0; x_true[17] = -2.0
y_obs = A @ x_true

# Basis pursuit
x_rec = basis_pursuit(A, y_obs)
err = np.linalg.norm(x_rec - x_true)
print(f"\nBasis pursuit recovery error: {err:.2e}")
print(f"Recovered support: {np.where(np.abs(x_rec) > 1e-3)[0].tolist()}"
      f"  (true: [5, 17])")
print("(Recovery succeeds even though RIP condition is not met --"
      " the condition is sufficient, not necessary.)")
\end{lstlisting}

\section{Chapter Summary}
\label{sec:summary-ch12}

\begin{itemize}
  \item \textbf{Sparsity in the frame setting.} When $m>n$, every
        signal $x$ has infinitely many representations $c$ with $Fc=x$;
        they form an affine family $c+\ker F$
        (Theorem~5.4.1). \emph{Basis pursuit} --- $\ell^1$ minimization
        over this family --- selects the sparsest, as
        Example~\ref{ex:triangular-sparsity} illustrates explicitly:
        $\ell^1$ finds the unique 1-sparse expansion of $x=f_1$, with
        $\ell^1$ norm $1$ vs.\ the canonical vector's $\tfrac43$.

  \item \textbf{The compressed sensing model.} A sensing matrix
        $A\in\R^{n\times m}$ ($n\ll m$) acquires $s$-sparse signals
        via $y=Ax$. Basis pursuit \eqref{eq:cs-bp} recovers $x$ exactly
        from $y$ under suitable conditions on $A$.

  \item \textbf{Coherence condition.} If the column coherence
        $\mu(A)$ satisfies $s<\tfrac12(1+1/\mu)$
        (Theorem~\ref{thm:coherence-recovery}), basis pursuit is
        guaranteed to recover any $s$-sparse signal. The coherence
        $\mu$ is exactly the quantity studied in Chapter~9 (Definition
        9.2.1), and ETFs minimize it (Proposition~\ref{prop:etf-optimal-sensing}).

  \item \textbf{Restricted isometry property.} The RIP
        (Definition~\ref{def:rip}) is a \emph{restricted frame
        condition}: $A$ satisfies RIP-$s$ iff every $s$-column
        subframe is a $(1\pm\delta_s)$-frame for $\R^n$
        (Proposition~\ref{prop:sensing-frame}). If $\delta_{2s}
        <\sqrt{2}-1$, exact recovery holds (Theorem~\ref{thm:rip-recovery}).
        Coherence controls RIP via $\delta_s\le(s-1)\mu$
        (Proposition~\ref{prop:coherence-rip}).

  \item \textbf{Frames are sensing matrices.} Every full-rank
        $n\times m$ matrix is a frame for $\R^n$ (Theorem~2.2.1).
        The entire frame-theoretic toolkit --- bounds, tightness,
        coherence, the Welch bound --- translates directly into
        properties of sensing matrices. ETFs are optimal sensing
        matrices in the coherence sense.
\end{itemize}

\section{Lab Exercises}
\label{sec:lab-ch12}

\begin{exercisebox}[title={Lab Exercise 12.1: Reproducing Example 12.1}]
Implement the $\ell^1$ minimization from
Section~\ref{sec:numpy-ch12} for the triangular frame. Verify
numerically that:
\begin{enumerate}[label=(\alph*)]
  \item The optimal $t^* = 1/\sqrt{3}$, giving $c^*=(1,0,0)^T$.
  \item For $x = \alpha f_2$ (any scalar multiple of the second frame
        vector), $\ell^1$ minimization finds the 1-sparse representation
        $c^* = (0,\alpha,0)^T$.
  \item For a \emph{generic} signal (not a multiple of any single frame
        vector), the $\ell^1$ minimizer has exactly two nonzero entries.
        (Try $x=(2,-1)^T$ and report the minimizing $c^*$.)
\end{enumerate}
\end{exercisebox}

\begin{exercisebox}[title={Lab Exercise 12.2: Basis Pursuit via LP}]
Implement \texttt{basis\_pursuit(A, y)} from
Section~\ref{sec:numpy-ch12} and test it on the following:
\begin{enumerate}[label=(\alph*)]
  \item The triangular frame with $x=f_1$: verify you get the same
        $c^*=(1,0,0)^T$ as Example~12.1.
  \item A random Gaussian sensing matrix ($n=20$, $m=100$) with a
        3-sparse true signal. Does basis pursuit recover exactly?
  \item Repeat (b) for sparsity levels $s=1,2,3,\ldots,10$. For which
        values of $s$ does recovery fail? Does the coherence bound
        (Theorem~\ref{thm:coherence-recovery}) correctly predict
        the failure threshold?
\end{enumerate}
\end{exercisebox}

\begin{exercisebox}[title={Lab Exercise 12.3: The Coherence--RIP Connection}]
For a random Gaussian sensing matrix with $n=20$ columns normalized to
unit norm and $m=80$:
\begin{enumerate}[label=(\alph*)]
  \item Compute the coherence $\mu$ and the estimated RIP-2 constant
        $\delta_2$ (using \texttt{rip\_constant} from
        Section~\ref{sec:numpy-ch12}).
  \item Verify the bound $\delta_2\le(2-1)\mu=\mu$
        (Proposition~\ref{prop:coherence-rip}).
  \item Repeat for $s=3,4,5$ to verify $\delta_s\le(s-1)\mu$.
  \item Is the bound tight (close to equality) or loose? What does this
        suggest about the relative sharpness of coherence-based vs.\
        RIP-based guarantees?
\end{enumerate}
\end{exercisebox}

\begin{exercisebox}[title={Lab Exercise 12.4: ETFs as Sensing Matrices}]
Compare the triangular frame ($n=2$, $m=3$, coherence $\mu=0.5$) to
two other sensing matrices of the same size:
\begin{enumerate}[label=(\alph*)]
  \item A random unit-column-norm $2\times3$ matrix (random Gaussian,
        columns normalized).
  \item The $2\times3$ matrix $A=[[1,0,1],[0,1,0]]$ (not normalized;
        normalize columns first).
\end{enumerate}
For each, compute $\mu$, the coherence bound on sparsity recovery
($s<(1+1/\mu)/2$), and $\delta_2$. Rank the three matrices by their
sensing quality (smallest $\mu$). Is the triangular frame (the ETF for
$n=2$, $m=3$) the best? If so, why --- connect to the Welch bound from
Chapter~9.
\end{exercisebox}

\begin{exercisebox}[title={Lab Exercise 12.5 (Challenge): Noisy Measurements}]
In practice, measurements are noisy: $y = Ax + e$ with $\norm{e}\le\eta$.
The noisy recovery program (LASSO or basis pursuit denoising) is
\[
  \hat x = \argmin_{z}\norm{z}_1
  \quad\text{subject to}\quad\norm{Az-y}\le\eta.
\]
If $A$ satisfies RIP-$2s$ with $\delta_{2s}<\sqrt{2}-1$, the recovery
error satisfies $\norm{\hat x - x}\le C\eta$ for a constant $C$
depending on $\delta_{2s}$.
\begin{enumerate}[label=(\alph*)]
  \item Implement noisy basis pursuit using \texttt{cvxpy} (install
        via \texttt{pip install cvxpy}): the constraint becomes
        \texttt{cp.norm(A @ z - y, 2) <= eta}.
  \item For a random $20\times80$ sensing matrix and a 2-sparse true
        signal, add Gaussian noise with $\sigma=0.01,0.05,0.1,0.2$
        (set $\eta=2\sigma\sqrt{n}$ as a rough bound on $\norm{e}$).
        Plot the recovery error $\norm{\hat x - x}$ as a function of
        $\sigma$.
  \item Verify empirically that the error scales roughly linearly
        with $\sigma$ (i.e.\ $\norm{\hat x-x}\approx C\sigma$), as the
        theory predicts.
\end{enumerate}
\end{exercisebox}

\newpage
\section{Supplementary Exercises}
\label{sec:extra-ch12}

\subsection*{Theoretical Exercises}
\begin{theoexercises}
	\item Recall the \emph{spark} of a matrix $A$ (Chapter~10, Remark
	following Definition~10.4.1): the smallest number of linearly
	dependent columns. Prove that if $x$ is $s$-sparse, $z$ is
	$s'$-sparse, $Ax=Az$, and $s+s'<\operatorname{spark}(A)$, then
	$x=z$. Conclude that a full-spark sensing matrix (spark $=n+1$,
	the maximum possible) gives the best possible uniqueness
	guarantee, recovering any $s$-sparse signal uniquely whenever
	$s\le\lfloor n/2\rfloor$.
	
	\item Prove that the $\ell^0$ ``norm'' $\norm{\cdot}_0$ satisfies the
	triangle inequality but \emph{fails} absolute homogeneity
	($\norm{cx}_0\ne|c|\norm{x}_0$ in general), and is therefore not
	a genuine norm despite the name.
	
	\item Explain why the coherence-based recovery threshold of
	Theorem~12.4.1, $s<\tfrac12(1+1/\mu)$, becomes vacuous (tends to
	infinity) as $\mu\to0$, and why this limiting case is consistent
	with, but not practically meaningful for, compressed sensing
	(recall $\mu=0$ forces $m\le n$, i.e.\ no redundancy to exploit).
	
	\item Prove that the RIP constant $\delta_s$ is non-decreasing in $s$:
	$\delta_s\le\delta_{s+1}$ for all $s$. (Hint: use the Cauchy
	interlacing theorem on the Gram matrix of an $(s+1)$-subset that
	extends a worst-case $s$-subset.)
	
	\item Using Proposition~12.4.4 ($\delta_s\le(s-1)\mu$) together with the
	RIP recovery condition $\delta_{2s}<\sqrt2-1$, derive a
	coherence-sufficient condition for RIP-based recovery, and show
	it gives a strictly \emph{more conservative} (smaller) guaranteed
	sparsity level than the direct coherence bound of Theorem~12.4.1,
	confirming the chapter's remark that coherence-based RIP bounds
	are pessimistic.
\end{theoexercises}

\subsection*{Computational Exercises}
\begin{compexercises}
	\item Compute the spark of the triangular frame's synthesis matrix by
	brute-force search over all column subsets, confirm it equals
	$3$, and verify T1's uniqueness guarantee directly: check that no
	other $1$-sparse vector produces the same measurement as a given
	$1$-sparse $x$.
	
	\item For a random $5\times10$ sensing matrix (small enough to check
	\emph{all} subsets exhaustively, not a random sample), compute
	the exact RIP constants for $s=1,\ldots,6$ and confirm they are
	non-decreasing, verifying T4. (Caution: using the subsampled
	\texttt{rip\_constant} from Lab Exercise~12.3 with a small
	\texttt{max\_subsets} can spuriously \emph{violate} monotonicity
	due to sampling noise --- use exhaustive
	\texttt{itertools.combinations} for this exercise.)
	
	\item For $\mu=0.05,0.1,0.2,0.3$, plot both the direct coherence
	threshold from Theorem~12.4.1 and the coherence-via-RIP threshold
	from T5 as functions of $\mu$, confirming the direct threshold is
	always larger (more permissive).
	
	\item Construct a random $10\times40$ sensing matrix, and for
	increasing sparsity levels $s=1,\ldots,8$, run basis pursuit on a
	random $s$-sparse signal and record whether recovery succeeds.
	Compare the empirical failure threshold to both the direct
	coherence bound (Theorem~12.4.1) and the spark-based uniqueness
	bound (T1, using $2s<\operatorname{spark}(A)$, estimated via
	partial search since exact spark computation is expensive for
	$m=40$).
\end{compexercises}
